\section{Conclusion}\label{sec:conclusion}

We characterized IRT person-fit statistics for LLM contamination detection across three 7--8B model families, establishing reliable SFT-stage detection while identifying boundary conditions: $\theta$-shift absorption necessitating $\theta$-anchored estimation for high-ability models, a detection-diagnosis tradeoff, and non-monotonic dose-response as a $\theta$-estimation artifact.
For practitioners, we recommend a two-stage workflow: accuracy $z$-scores for efficient screening, followed by $\ell_z^*$ analysis when mechanistic understanding is needed---distinguishing \emph{which} items drive aberrance rather than merely flagging contamination.
Beyond detection, person-fit statistics capture response-pattern anomalies that persist after RL post-training erases memorization-based signals, complementing logit-based methods.
Future work should address pretraining-stage signatures, broader model scales, and cross-benchmark generalization with larger item banks.
